% PDF p. 316
\AppendixSection{6A}{体积展开与正规坐标}
\AppendixHeading{6A1}{正规坐标的概念和定义}
设 $(M,g)$ 为 $n$ 维黎曼流形，$p\in M$ 为一个固定点。对于坐标系 $(U,\varphi)$，其中 $\varphi\colon U\to\R^n$ 是满足 $\varphi(p)=0$ 的微分同胚。若它满足下列性质，就称它是以 $p$ 为中心的正规坐标系。
\begin{enumerate}
\item \textbf{将测地线拉直。}
在这组坐标下，从 $p$ 出发的每条测地线都表示为过原点的直线。更确切地说，若 $v\in T_pM$ 是一个切向量，则测地线 $\gamma_v(t)=\exp_p(tv)$ 在 $0$ 附近一个足够小的区间内满足 $\varphi(\gamma_v(t))=tv$。

这表明指数映射 $\exp_p\colon T_pM\to M$ 给出坐标卡，其坐标映射为 $\varphi(x)=\log_p(x)$。这里的 $\log_p$ 是 $\exp_p$ 在 $p$ 附近的逆。
\item \textbf{原点处的度量为欧氏度量。}
在以 $p$ 为中心的正规坐标中，黎曼度量 $g$ 在 $p$ 处的分量取标准欧氏形式
\[
  g_{ij}(0)=\delta_{ij},
\]
其中 $\delta_{ij}$ 是克罗内克符号（Kronecker delta）。
\item \textbf{克里斯托费尔符号消失。}
克里斯托费尔符号 $\Gamma^k_{ij}$ 在 $p$ 处为零，
\[
  \Gamma^k_{ij}(0)=0,
\]
这意味着度量分量的一阶导数在 $p$ 处也为零。
\end{enumerate}
综上，利用指数映射把 $M$ 上 $p$ 的邻域与 $T_pM\simeq\R^n$ 中 0 的邻域等同，就能构造 $p$ 处的正规坐标。在这组坐标中，流形在 $p$ 处“看起来是平的”。

\AppendixSubheading{6A1.1}{正规坐标举例}
\begin{enumerate}
\item \textbf{欧氏空间 $\R^n$。}
在配备标准欧氏度量的 $\R^n$ 中，通常的笛卡儿坐标在每一点处已经是正规坐标。对任意 $p\in\R^n$，若将 $p$ 平移至原点，度量仍为标准欧氏度量，测地线为直线，克里斯托费尔符号恒为零。
\item \textbf{二维球面 $S^2$。}
考虑配有标准球面度量的单位球面 $S^2\subset\R^3$。令 $p$ 为北极，例如 $p=(0,0,1)$。可以用指数映射定义 $p$ 处的正规坐标。
\begin{itemize}
\item 将 $T_pS^2$ 与 $\R^3$ 中的平面 $z=0$ 等同。
% PDF p. 317
\item 对于 $T_pS^2$ 中的切向量 $v$，指数映射 $\exp_p(v)$ 给出 $S^2$ 上的一个点。该点沿以 $p$ 为起点、以 $v$ 为初速度的唯一大圆到达。

在这种情形下，若把切平面中的 $v$ 写为 $(v_1,v_2)$，则球面上的点在正规坐标中表示为 $v=\exp_p(v)$。%
\translatornote{原书左端印为 $v$。此处应理解为球面点 $\exp_p(v)$ 的正规坐标是 $v=(v_1,v_2)$。}%
一个显式表达式为（适用于 $\lVert v\rVert$ 较小时）
\[
  \exp_p(v)=\left(
    \frac{\sin(\lVert v\rVert)}{\lVert v\rVert}v_1,
    \frac{\sin(\lVert v\rVert)}{\lVert v\rVert}v_2,
    \cos(\lVert v\rVert)\right).
\]
这里，$(v_1,v_2)$ 作为 $S^2$ 上 $p$ 附近的正规坐标。在 $p$ 附近，用这些坐标表示的度量与平坦度量的差从二阶项开始。
\end{itemize}
\item \textbf{双曲平面 $\mathbb H^2$。}
在双曲平面上，也能用指数映射引入以点 $p$ 为中心的正规坐标。若采用庞加莱圆盘或上半平面模型来表示 $\mathbb H^2$，则通过指数映射将切空间 $T_p\mathbb H^2$（同构于 $\R^2$）映入 $\mathbb H^2$，可得到 $p$ 附近的正规坐标。在这组坐标下，经过 $p$ 的测地线表现为从 $T_p\mathbb H^2$ 原点出发的直线，$p$ 处的度量为欧氏度量，曲率的影响只出现在更高阶项中。
\end{enumerate}

\AppendixHeading{6A2}{仿射联络}
\AppendixSubheading{6A2.1}{引言}
光滑流形 $M$ 上的仿射联络是黎曼几何的重要工具。它使我们可以比较 $M$ 上不同点处的向量，定义如何沿曲线对向量场求导，用光滑且一致的方式沿曲线移动向量，并在流形上构造测地线。这些概念是理解内在几何及其与曲率和挠率关系的基础。

\AppendixSubheading{6A2.2}{方向导数}
回顾偏导数 $f_x(x,y)$ 和 $f_y(x,y)$。它们分别表示固定 $y$ 而改变 $x$、固定 $x$ 而改变 $y$ 时 $f$ 的变化率。现在需要讨论 $x$ 和 $y$ 同时变化时，如何求 $f$ 的变化率（Paul's Online Notes 2025）。

先定义标量函数 $f$ 的方向导数。设 $X=(x_1,x_2,x_3)$ 且 $f(X)=f(x_1,x_2,x_3)$。标量函数的方向导数定义为
\begin{equation*}
  D_X f=\lim_{t\to0}
    \frac{f(p_1+tx_1,p_2+tx_2,p_3+tx_3)-f(p_1,p_2,p_3)}{t},
  \tag{6A1}\label{eq:6A1}
\end{equation*}
\begin{equation*}
  =\langle\nabla_p f,X\rangle,
  \tag{6A2}\label{eq:6A2}
\end{equation*}
% PDF p. 318
其中函数 $f$ 的梯度为
\[
  \nabla_p f=\left.
  \begin{bmatrix}
  \dfrac{\partial f}{\partial x_1}\\[4pt]
  \dfrac{\partial f}{\partial x_2}\\[4pt]
  \dfrac{\partial f}{\partial x_3}
  \end{bmatrix}\right|_{X=p}.
\]

接着定义向量的方向导数。设 $Y\colon\R^3\to\R^3$ 是向量场，$X\in\R^3$ 是一个向量（方向）。定义雅可比矩阵
\[
  J=\frac{\partial Y}{\partial X^T}=
    \begin{bmatrix}
      \dfrac{\partial y_1}{\partial x_1}&
      \dfrac{\partial y_1}{\partial x_2}&
      \dfrac{\partial y_1}{\partial x_3}\\[4pt]
      \dfrac{\partial y_2}{\partial x_1}&
      \dfrac{\partial y_2}{\partial x_2}&
      \dfrac{\partial y_2}{\partial x_3}\\[4pt]
      \dfrac{\partial y_3}{\partial x_1}&
      \dfrac{\partial y_3}{\partial x_2}&
      \dfrac{\partial y_3}{\partial x_3}
    \end{bmatrix}.
\]
那么 $Y$ 在点 $p\in\R^3$ 处沿方向 $X$ 的方向导数为
\begin{equation*}
\begin{aligned}
  D_XY(p)
    &=\left.\frac{d}{dt}Y(p+tX)\right|_{t=0},\\
    &=\begin{bmatrix}
      \dfrac{\partial y_1}{\partial x^1}&
      \dfrac{\partial y_1}{\partial x^2}&
      \dfrac{\partial y_1}{\partial x^3}\\[4pt]
      \dfrac{\partial y_2}{\partial x^1}&
      \dfrac{\partial y_2}{\partial x^2}&
      \dfrac{\partial y_2}{\partial x^3}\\[4pt]
      \dfrac{\partial y_3}{\partial x^1}&
      \dfrac{\partial y_3}{\partial x^2}&
      \dfrac{\partial y_3}{\partial x^3}
    \end{bmatrix}
    \begin{bmatrix}x^1\\x^2\\x^3\end{bmatrix},\\
    &=JX.
\end{aligned}
\tag{6A3}\label{eq:6A3}
\end{equation*}
该导数描述的是沿 $X$ 方向的直线上向量场 $Y$ 的变化，未计入曲率。

\AppendixSubheading{6A2.3}{坐标表示}
定义
\begin{align*}
  X&=X^i\frac{\partial}{\partial x^i}
     =\sum_i X^i\frac{\partial}{\partial x^i},\\
  Y&=Y^j(x)\frac{\partial}{\partial x^j}
     =\sum_j Y^j(x)\frac{\partial}{\partial x^j}.
\end{align*}

% PDF p. 319
由式 \eqref{eq:6A3} 可得
\[
  D_XY(p)=
  \begin{bmatrix}
    \displaystyle\sum_i x^i\frac{\partial Y^1}{\partial x^i}\\[5pt]
    \displaystyle\sum_i x^i\frac{\partial Y^2}{\partial x^i}\\[5pt]
    \displaystyle\sum_i x^i\frac{\partial Y^3}{\partial x^i}
  \end{bmatrix}.
\]
也可以写成
\begin{equation*}
  D_XY(p)=\sum_j\sum_i X^i\frac{\partial Y^j}{\partial x^i}\frac{\partial}{\partial x^j}
    =X^i\frac{\partial Y^j}{\partial x^i}\frac{\partial}{\partial x^j}.
  \tag{6A4}\label{eq:6A4}
\end{equation*}
这正是协变导数公式中的第一项
\[
  \nabla_XY=D_XY+(\text{联络修正项}).
\]
所以，在克里斯托费尔符号为零的平坦空间 $\R^n$ 中，$\nabla_XY=D_XY$。

\begin{statement}{例 1A.1}{球面上的向量场}
定义向量场 $Y(\varnothing)=[-\sin(\varnothing),\cos(\varnothing),0]$。在笛卡儿标架中，它是指向 $\varnothing$ 增大方向的向量场。

\textbf{用坐标表示 $Y$。}
\[
  Y=Y^j(x)\frac{\partial}{\partial x^j}
   =(-\sin\varnothing)\frac{\partial}{\partial x}
    +(\cos\varnothing)\frac{\partial}{\partial y}
    +0\frac{\partial}{\partial z}.
\]
因此，分量为 $Y^1=-\sin\varnothing$、$Y^2=\cos\varnothing$、$Y^3=0$。

\textbf{用坐标表示 $X$。}
设 $X$ 是指向经度 $\varnothing$ 增大方向的向量场。在固定 $\theta$ 时，单位球面上相应的笛卡儿坐标表示为%
\translatornote{原书此式将经度坐标向量与单位切向量等同。按下文球面度量 $g_{\varnothing\varnothing}=\sin^2\theta$，等式右端为 $(\sin\theta)^{-1}\partial_\varnothing$（在两极之外），在赤道上才与 $\partial_\varnothing$ 相等。以下取 $X^\varnothing=1$，即采用坐标向量。}%
\[
  X=\frac{\partial}{\partial\varnothing}
    =-\sin\varnothing\frac{\partial}{\partial x}
     +\cos\varnothing\frac{\partial}{\partial y}
     +0\frac{\partial}{\partial z},
\]
所以
\[
  X^1=-\sin\varnothing,\quad X^2=\cos\varnothing,\quad X^3=0.
\]

\textbf{计算方向导数 $D_XY$。}
\[
  D_XY(p)=\sum_j\sum_i X^i\frac{\partial Y^j}{\partial x^i}\frac{\partial}{\partial x^j}
     =X^i\frac{\partial Y^j}{\partial x^i}\frac{\partial}{\partial x^j}.
\]
由于 $Y^j$ 只依赖 $\varnothing$，而且 $X^i$ 和 $Y^j$ 都是 $\varnothing$ 的函数，需要计算导数 $\partial Y^j/\partial x^i$。但是在笛卡儿坐标中，$\varnothing$ 不是 $x,y,z$ 的线性函数，所以计算较麻烦。

% PDF p. 320
由于所有量只随 $\varnothing$ 变化，可用链式法则对 $\varnothing$ 求导
\begin{equation*}
  \frac{dY}{d\varnothing}
   =\begin{bmatrix}-\cos\varnothing\\-\sin\varnothing\\0\end{bmatrix}.
  \tag{6A5}\label{eq:6A5}
\end{equation*}
因此，式 \eqref{eq:6A4} 化为
\begin{equation*}
  D_XY=\left(X^\varnothing\frac{dY^j}{d\varnothing}\right)
      \frac{\partial}{\partial x^j}.
  \tag{6A6}\label{eq:6A6}
\end{equation*}
注意
\begin{equation*}
  X^\varnothing=1.
  \tag{6A7}\label{eq:6A7}
\end{equation*}
把式 \eqref{eq:6A5} 和式 \eqref{eq:6A7} 代入式 \eqref{eq:6A6}，得
\begin{equation*}
  D_XY=-\cos\varnothing\frac{\partial}{\partial x}
        -\sin\varnothing\frac{\partial}{\partial y}
        +0\frac{\partial}{\partial z}.
  \tag{6A8}\label{eq:6A8}
\end{equation*}
\end{statement}

\AppendixSubheading{6A2.4}{仿射联络的定义}
光滑流形 $M$ 上的仿射联络使我们可以
\begin{itemize}
\item 比较 $M$ 上不同点处的向量。
\item 定义如何沿曲线对向量场求导。
\item 用光滑而一致的方式沿曲线移动向量。
\end{itemize}
形式上，仿射联络 $\nabla$ 是一个双线性映射
\begin{equation*}
  \nabla\colon\mathfrak X(M)\times\mathfrak X(M)\to\mathfrak X(M),
  \qquad (X,Y)\mapsto\nabla_XY,
  \tag{6A9}\label{eq:6A9}
\end{equation*}
其中 $\mathfrak X(M)$ 是 $M$ 上光滑向量场的空间。它满足下列性质。
\begin{itemize}
\item \textbf{关于 $X$ 和 $Y$ 的线性性。}
\begin{equation*}
  \nabla_{fX_1+gX_2}Y=f\nabla_{X_1}Y+g\nabla_{X_2}Y,
  \tag{6A10}\label{eq:6A10}
\end{equation*}
\begin{equation*}
  \nabla_X(aY_1+bY_2)=a\nabla_XY_1+b\nabla_XY_2.
  \tag{6A11}\label{eq:6A11}
\end{equation*}
\item \textbf{莱布尼茨法则。} 对光滑函数 $f$，
\begin{equation*}
  \nabla_X(fY)=f\nabla_XY+X(f)Y.
  \tag{6A12}\label{eq:6A12}
\end{equation*}
\end{itemize}

为什么这个定义可以实现直观上所期待的目标，也就是比较向量、对向量求导以及移动向量？下面逐一说明仿射联络的定义如何做到这一点。

% PDF p. 321
\AppendixSubheading{6A2.5}{比较不同点处的向量}
不同点处的向量属于不同的切空间。例如 $v\in T_pM$，而 $w_q\in T_qM$。要比较这两个向量，需要额外的结构。

有了仿射联络，可以定义一个称为平行移动（parallel transport）的过程。
\begin{itemize}
\item 如果沿曲线 $\gamma(t)$ 从 $p=\gamma(0)$ 移到 $q=\gamma(1)$，联络规定如何沿该曲线一致地“拖动”向量。
\item 根据这个规则，可将 $v_p\in T_pM$ 移动为一个新向量 $v_q\in T_qM$，从而进行比较。
\item 平行移动由下式定义
\[
  \nabla_{\dot\gamma(t)}V(t)=0,
\]
其中 $V(t)$ 是沿曲线的向量场。

图 \ref{fig:6.A1} 解释了为何 $\nabla_{\dot\gamma(t)}V(t)=0$ 定义了平行移动。图中的黑色曲线是路径 $\gamma(t)$，这里取一条半圆。灰色箭头是切向量 $\dot\gamma(t)$，表示运动方向。深红色箭头表示沿曲线平行移动的向量场 $V(t)$。图 \ref{fig:6.A1} 说明了这个方程为何合理。条件 $\nabla_{\dot\gamma(t)}V(t)=0$ 表示，沿曲线 $\gamma(t)$ 运动时，用流形的几何来度量，向量场 $V(t)$ 沿运动方向保持不变。
\item 它保证 $V(t)$ 始终属于每一点处的切空间。
\item 向量的长度及其相对于联络的方向保持不变。%
\translatornote{联络与度量相容时，平行移动才保证长度保持不变。图中的红色箭头可理解为在欧氏平面内沿半圆路径平行移动的向量。它们属于平面的切空间。}
\item 这是沿曲线光滑且一致地移动向量的方法。
\end{itemize}

\begin{figure}[tbp]
\centering
% PDF p. 321. Replace the three legend labels within the original plot.
\begin{tikzpicture}
\node[inner sep=0pt,outer sep=0pt,anchor=south west] (plot) at (0,0)
  {\includegraphics[width=.86\linewidth]{assets/ch06/fig06_a01.pdf}};
\begin{scope}[x={(plot.south east)},y={(plot.north west)}]
\fill[white] (.461,.151) rectangle (.661,.317);
\node[anchor=west,inner sep=0pt,font=\scriptsize] at (.463,.295) {曲线 $\gamma(t)$};
\node[anchor=west,inner sep=0pt,font=\scriptsize] at (.463,.239) {切向量 $\dot\gamma(t)$};
\node[anchor=west,inner sep=0pt,font=\scriptsize] at (.463,.184) {平行移动 $V(t)$};
\end{scope}
\end{tikzpicture}

\FigureCaption{6.A1}{平行移动为何如此定义。}
\end{figure}

\AppendixSubheading{6A2.6}{联络定义协变导数}
协变导数的概念由偏导数的概念发展而来，是测地线概念的基础（Randy 2023）。在数学中，协变导数是规定沿流形切向量如何求导的方法（Wikipedia）。
% PDF p. 322
向量场作用于标量场时相当于一个导数，但还需定义一个向量场在另一个向量场上的作用。本节引入协变导数 $\nabla_V$，借助它，向量场 $V$ 可以类似于求导那样作用于其他向量场。

\Subheading{6A2.6.1\quad 协变导数的定义}
设 $M$ 是光滑流形（例如曲面），$\nabla$ 是仿射联络（例如黎曼流形上的列维-奇维塔联络）。给定
\begin{itemize}
\item $M$ 上的向量场 $Y$。
\item $M$ 上的方向向量场 $X$（通常指各向量都具有单位长度的向量场，即只表示方向）。
\end{itemize}
$Y$ 沿 $X$ 方向的协变导数记为 $\nabla_XY$，它说明在考虑流形曲率的情况下 $Y$ 沿 $X$ 如何变化。

更确切地说，协变导数 $\nabla_XY$ 定义为%
\translatornote{这一投影表达式适用于带诱导度量的欧氏子流形的列维-奇维塔联络。一般仿射联络由式 \eqref{eq:6A9}–\eqref{eq:6A12} 的性质刻画。}%
\[
  \nabla_XY=\text{将普通导数 $D_XY$ 投影到 $M$ 的切空间所得的向量。}
\]

\Subheading{6A2.6.2\quad 协变导数的主要性质}
设 $f$ 是光滑标量函数，$X,Y,Z$ 是向量场。
\begin{itemize}
\item \textbf{关于 $X$ 和 $Y$ 的线性性。}
\begin{equation*}
  \nabla_{aX+bY}Z=a\nabla_XZ+b\nabla_YZ,
  \tag{6A13}\label{eq:6A13}
\end{equation*}
\item \textbf{莱布尼茨法则。}
\begin{equation*}
  \nabla_X(fY)=X(f)Y+f\nabla_XY,
  \tag{6A14}\label{eq:6A14}
\end{equation*}
\item 在黎曼流形上，列维-奇维塔联络满足

度量相容性给出 $X\langle Y,Z\rangle=\langle\nabla_XY,Z\rangle+\langle Y,\nabla_XZ\rangle$。

无挠性给出 $\nabla_XY-\nabla_YX=[X,Y]$。
\end{itemize}

\Subheading{6A2.6.3\quad 协变导数的一般公式}
设 $M$ 是坐标为 $(x^1,\dots,x^n)$ 的光滑流形，并设
\[
  X=X^i\frac{\partial}{\partial x^i},\qquad
  Y=Y^j\frac{\partial}{\partial x^j}
\]
是 $M$ 上的向量场。协变导数 $\nabla_XY$ 定义为
\begin{equation*}
  \nabla_XY=X^i\left(\frac{\partial Y^j}{\partial x^i}+\Gamma^j_{ik}Y^k\right)
    \frac{\partial}{\partial x^j},
  \tag{6A15}\label{eq:6A15}
\end{equation*}

% PDF p. 323
其中
\begin{itemize}
\item $\dfrac{\partial Y^j}{\partial x^i}$ 为通常的偏导数。
\item $\Gamma^j_{ik}$ 是克里斯托费尔符号（联络系数），满足
\begin{equation*}
  \nabla_{\partial_i}\partial_k=\Gamma^j_{ik}\partial_j.
  \tag{6A16}\label{eq:6A16}
\end{equation*}
\item 重复指标表示求和。
\end{itemize}

\noindent\textbf{证明}\par
令
\begin{equation*}
  X=X^i\partial_i=X^i\frac{\partial}{\partial x^i},\qquad
  Y=Y^k\partial_k=Y^k\frac{\partial}{\partial x^k}.
  \tag{6A17}\label{eq:6A17}
\end{equation*}
则
\begin{equation*}
  \nabla_XY=\nabla_{X^i\partial_i}(Y^k\partial_k).
  \tag{6A18}\label{eq:6A18}
\end{equation*}
根据关于 $X$ 的线性性（式 \eqref{eq:6A13}），
\begin{equation*}
  \nabla_XY=X^i\nabla_{\partial_i}(Y^k\partial_k).
  \tag{6A19}\label{eq:6A19}
\end{equation*}
现在应用莱布尼茨法则 \eqref{eq:6A14}，得
\begin{equation*}
  \nabla_{\partial_i}(Y^k\partial_k)=\partial_i(Y^k)\partial_k
      +Y^k\nabla_{\partial_i}\partial_k.
  \tag{6A20}\label{eq:6A20}
\end{equation*}
将式 \eqref{eq:6A16} 代入式 \eqref{eq:6A20}，得
\[
  \nabla_{\partial_i}(Y^k\partial_k)=\partial_i(Y^k)\partial_k
    +Y^k\Gamma^j_{ik}\partial_j,
\]
也就是
\begin{equation*}
  \nabla_{\partial_i}(Y^k\partial_k)=\frac{\partial Y^k}{\partial x^i}\partial_k
    +Y^k\Gamma^j_{ik}\partial_j.
  \tag{6A21}\label{eq:6A21}
\end{equation*}

接下来重排并重命名哑指标。
\begin{itemize}
\item 将式 \eqref{eq:6A21} 第一项中的 $k$ 改记为 $j$，便得到 $\dfrac{\partial Y^j}{\partial x^i}\partial_j$。于是
\begin{equation*}
  \nabla_{\partial_i}Y
  =\nabla_{\partial_i}(Y^k\partial_k)
  =\frac{\partial Y^j}{\partial x^i}\partial_j+Y^k\Gamma^j_{ik}\partial_j
  =\left(\frac{\partial Y^j}{\partial x^i}+Y^k\Gamma^j_{ik}\right)\partial_j.
  \tag{6A22}\label{eq:6A22}
\end{equation*}
\end{itemize}
把式 \eqref{eq:6A22} 代入式 \eqref{eq:6A19}，得
\[
  \nabla_XY=X^i\nabla_{\partial_i}(Y^k\partial_k)
    =X^i\left(\frac{\partial Y^j}{\partial x^i}+\Gamma^j_{ik}Y^k\right)
     \frac{\partial}{\partial x^j}.
\]
% PDF p. 324
证明完毕。

\Subheading{6A2.6.4\quad 克里斯托费尔符号公式（列维-奇维塔联络）}
给定黎曼度量 $g_{ij}$，第二类克里斯托费尔符号为
\begin{equation*}
  \Gamma^k_{ij}
   =\frac12 g^{kl}\left(\frac{\partial g_{jl}}{\partial x^i}
       +\frac{\partial g_{il}}{\partial x^j}
       -\frac{\partial g_{ij}}{\partial x^l}\right),
  \tag{6A23}\label{eq:6A23}
\end{equation*}
其中 $g^{kl}$ 是逆度量 $g^{-1}$ 的分量。对重复指标 $l$（$l=1,\dots,n$）求和，其中 $n$ 是流形维数，$x^i$ 是坐标。

以下证明第二类（即带上指标的）克里斯托费尔符号的公式，假设所用的联络是列维-奇维塔联络。

\noindent\textbf{证明}

\noindent\textbf{假设}\par
列维-奇维塔联络由下面两个性质唯一确定。

\noindent\textbf{度量相容性}
\[
  \nabla g=0\qquad\text{（度量的协变导数为零）},
\]
\noindent\textbf{无挠性}
\[
  \nabla_XY-\nabla_YX=[X,Y]\qquad\text{（挠率为零）},
\]
其中 $[X,Y]$ 是李括号。下面将
\begin{enumerate}
\item 利用度量相容性，得到度量导数与克里斯托费尔符号之间的关系。
\item 利用无挠条件给出的对称性。
\item 用所得方程组求解克里斯托费尔符号。
\end{enumerate}
证明分为两步。

\noindent\textbf{步骤 1。从度量相容性出发。}\par
回顾度量相容性蕴含的等式
\begin{equation*}
  X\langle Y,Z\rangle=\langle\nabla_XY,Z\rangle+\langle Y,\nabla_XZ\rangle.
  \tag{6A24}\label{eq:6A24}
\end{equation*}
设 $\partial_i=\dfrac{\partial}{\partial x^i}$、$\partial_j=\dfrac{\partial}{\partial x^j}$、
$\partial_k=\dfrac{\partial}{\partial x^k}$ 是局部坐标向量场。对于任意这样的向量场 $\partial_i,\partial_j,\partial_k$，度量相容性给出
\begin{equation*}
  X\langle Y,Z\rangle=\frac{\partial}{\partial x^i}g(\partial_j,\partial_k)
    =\frac{\partial g_{jk}}{\partial x_i},
  \tag{6A25}\label{eq:6A25}
\end{equation*}
\begin{equation*}
  \langle\nabla_XY,Z\rangle=g(\nabla_{\partial_i}\partial_j,\partial_k),
  \tag{6A26}\label{eq:6A26}
\end{equation*}
% PDF p. 325
\begin{equation*}
  \langle Y,\nabla_XZ\rangle=g(\partial_j,\nabla_{\partial_i}\partial_k).
  \tag{6A27}\label{eq:6A27}
\end{equation*}
合并式 \eqref{eq:6A25} 至式 \eqref{eq:6A27}，得
\begin{equation*}
  \frac{\partial g_{jk}}{\partial x^i}
    =g(\nabla_{\partial_i}\partial_j,\partial_k)
     +g(\partial_j,\nabla_{\partial_i}\partial_k).
  \tag{6A28}\label{eq:6A28}
\end{equation*}
度量相容性也意味着
\begin{equation*}
  \nabla g=0.
  \tag{6A29}\label{eq:6A29}
\end{equation*}
因此，内积作为度量的协变导数为零，即
\begin{equation*}
  X(g(Y,Z))=g(\nabla_XY,Z)+g(Y,\nabla_XZ).
  \tag{6A30}\label{eq:6A30}
\end{equation*}
类似于式 \eqref{eq:6A24}，式 \eqref{eq:6A30} 导出
\begin{equation*}
  \frac{\partial g_{jk}}{\partial x^i}
    =g(\nabla_{\partial_i}\partial_j,\partial_k)
     +g(\partial_j,\nabla_{\partial_i}\partial_k).
  \tag{6A31}\label{eq:6A31}
\end{equation*}
由式 \eqref{eq:6A16}，
\begin{equation*}
  \nabla_{\partial_i}\partial_k=\Gamma^m_{ik}\partial_m,
  \tag{6A32}\label{eq:6A32}
\end{equation*}
\begin{equation*}
  \nabla_{\partial_i}\partial_j=\Gamma^m_{ij}\partial_m.
  \tag{6A33}\label{eq:6A33}
\end{equation*}
将式 \eqref{eq:6A32} 和式 \eqref{eq:6A33} 代入式 \eqref{eq:6A31}，得
\begin{equation*}
  \frac{\partial g_{jk}}{\partial x^i}
   =g(\Gamma^m_{ij}\partial_m,\partial_k)+g(\partial_j,\Gamma^m_{ik}\partial_m).
  \tag{6A34}\label{eq:6A34}
\end{equation*}
度量对两个变量均线性，因此可以提出系数
\begin{equation*}
  g(\Gamma^m_{ij}\partial_m,\partial_k)
   =\Gamma^m_{ij}g(\partial_m,\partial_k)=\Gamma^m_{ij}g_{mk},
  \tag{6A35}\label{eq:6A35}
\end{equation*}
\begin{equation*}
  g(\partial_j,\Gamma^m_{ik}\partial_m)
   =\Gamma^m_{ik}g(\partial_j,\partial_m)=\Gamma^m_{ik}g_{jm}.
  \tag{6A36}\label{eq:6A36}
\end{equation*}
将式 \eqref{eq:6A35} 和式 \eqref{eq:6A36} 代入式 \eqref{eq:6A34}，得
\begin{equation*}
  \frac{\partial g_{jk}}{\partial x^i}
   =\Gamma^m_{ij}g_{mk}+\Gamma^m_{ik}g_{jm}
   =g_{mk}\Gamma^m_{ij}+g_{mj}\Gamma^m_{ik}.
  \tag{6A37}\label{eq:6A37}
\end{equation*}
类似地，得到
\begin{equation*}
  \frac{\partial g_{ik}}{\partial x^j}=g_{mk}\Gamma^m_{ji}+g_{mi}\Gamma^m_{jk},
  \tag{6A38}\label{eq:6A38}
\end{equation*}
\begin{equation*}
  \frac{\partial g_{ij}}{\partial x^k}=g_{mj}\Gamma^m_{ki}+g_{mi}\Gamma^m_{kj}.
  \tag{6A39}\label{eq:6A39}
\end{equation*}

% PDF p. 326
由于 $\nabla$ 无挠，
\[
  \Gamma^m_{ij}=\Gamma^m_{ji}.
\]
由式 \eqref{eq:6A37} 至式 \eqref{eq:6A39} 得
\begin{equation*}
  2g_{mk}\Gamma^m_{ij}
   =\frac{\partial g_{jk}}{\partial x^i}+\frac{\partial g_{ik}}{\partial x^j}
    -\frac{\partial g_{ij}}{\partial x^k}.
  \tag{6A40}\label{eq:6A40}
\end{equation*}
现在用 $g^{mk}$ 乘式 \eqref{eq:6A40} 两端，解出 $\Gamma^m_{ij}$%
\translatornote{为区分自由指标与求和指标，可先将式 \eqref{eq:6A40} 左端的求和指标 $m$ 改为 $s$，再乘 $g^{mk}$，此时 $g^{mk}g_{sk}=\delta^m_s$。}%
\begin{equation*}
\begin{aligned}
  2\Gamma^m_{ij}
    &=g^{mk}\left(\frac{\partial g_{jk}}{\partial x^i}
       +\frac{\partial g_{ik}}{\partial x^j}-\frac{\partial g_{ij}}{\partial x^k}\right),
       \qquad\text{或写为}\\
  \Gamma^m_{ij}
    &=\frac12 g^{mk}\left(\frac{\partial g_{jk}}{\partial x^i}
       +\frac{\partial g_{ik}}{\partial x^j}-\frac{\partial g_{ij}}{\partial x^k}\right).
\end{aligned}
\tag{6A41}\label{eq:6A41}
\end{equation*}
将 $k$ 改记为 $l$、将 $m$ 改记为 $k$ 后，式 \eqref{eq:6A41} 化为
\begin{equation*}
  \Gamma^k_{ij}=\frac12 g^{kl}
     \left(\frac{\partial g_{jl}}{\partial x^i}
          +\frac{\partial g_{il}}{\partial x^j}-\frac{\partial g_{ij}}{\partial x^l}\right).
  \tag{6A42}\label{eq:6A42}
\end{equation*}
证明完毕。

\Subheading{6A2.6.5\quad 二维球面上的协变导数通式}
考虑坐标为 $(\theta,\varnothing)$ 的二维球面，其度量为
\begin{equation*}
  g=\begin{bmatrix}1&0\\0&\sin^2\theta\end{bmatrix}.
  \tag{6A43}\label{eq:6A43}
\end{equation*}
设向量场为
\begin{equation*}
  X=X^1(\theta,\varnothing)\partial_\theta+X^2(\theta,\varnothing)\partial_\varnothing,
  \tag{6A44}\label{eq:6A44}
\end{equation*}
\begin{equation*}
  Y=Y^1(\theta,\varnothing)\partial_\theta+Y^2(\theta,\varnothing)\partial_\varnothing.
  \tag{6A45}\label{eq:6A45}
\end{equation*}
先计算二维球面上的第二类克里斯托费尔符号。采用坐标 $x^1=\theta$、$x^2=\varnothing$，由式 \eqref{eq:6A43} 得
\begin{gather*}
  g_{11}=1,\quad g_{12}=g_{21}=0,\quad g_{22}=\sin^2\theta,\\
  g^{-1}=\begin{bmatrix}1&0\\0&\dfrac{1}{\sin^2\theta}\end{bmatrix},\\
  g^{11}=1,\quad g^{12}=g^{21}=0,\quad g^{22}=\frac{1}{\sin^2\theta}.
\end{gather*}
% PDF p. 327
对各个分量应用公式，得二维球面上的非零克里斯托费尔符号
\begin{align*}
  \Gamma^1_{22}
    &=\frac12g^{11}\left(\frac{\partial g_{21}}{\partial\theta}
        +\frac{\partial g_{21}}{\partial\varnothing}
        -\frac{\partial g_{22}}{\partial\theta}\right)
      =\frac12\left(-\frac{\partial(\sin^2\theta)}{\partial\theta}\right)
      =-\sin\theta\cos\theta,\\
  \Gamma^2_{12}&=\Gamma^2_{21},\\
  \Gamma^2_{12}&=\frac12g^{22}\frac{\partial g_{22}}{\partial\theta}
    =\frac12\frac{1}{\sin^2\theta}\,2\sin\theta\cos\theta
    =\frac{\cos\theta}{\sin\theta}=\cot\theta.
\end{align*}
由式 \eqref{eq:6A15} 得协变导数 $\nabla_XY$
\begin{align*}
  \nabla_XY^1&=X^1\frac{\partial Y^1}{\partial\theta}
      +X^2\left(\frac{\partial Y^1}{\partial\varnothing}
        -\sin\theta\cos\theta\,Y^2\right),\\
  \nabla_XY^2&=X^1\left(\frac{\partial Y^2}{\partial\theta}+\frac{Y^2}{\tan\theta}\right)
      +X^2\left(\frac{\partial Y^2}{\partial\varnothing}+\frac{Y^1}{\tan\theta}\right).
\end{align*}

\Subheading{6A2.6.6\quad 庞加莱球模型度量下双曲平面极坐标 $(r,\theta)$ 中协变导数的通式}
极坐标下的双曲度量为
\[
  g=\begin{bmatrix}1&0\\0&\sinh^2r\end{bmatrix}.
\]
于是
\begin{gather*}
  g_{11}=1,\quad g_{12}=g_{21}=0,\quad g_{22}=\sinh^2r,\\
  g^{-1}=\begin{bmatrix}1&0\\0&\dfrac{1}{\sinh^2r}\end{bmatrix},\\
  g^{11}=1,\quad g^{12}=g^{21}=0,\quad g^{22}=\frac{1}{\sinh^2r}.
\end{gather*}
定义向量场
\begin{align*}
  X&=X^1(r,\theta)\partial_r+X^2(r,\theta)\partial_\theta,\\
  Y&=Y(r,\theta)\partial_r+Y^2(r,\theta)\partial_\theta.
\end{align*}
再计算双曲平面上的非零克里斯托费尔符号%
\translatornote{原式中间的 $\sin,\cos$ 应分别为 $\sinh,\cosh$，末项应为 $1/\tanh r$。下一式使用的是 $1/\tanh r$。}%
\begin{align*}
 \Gamma^1_{22}
  &=\frac12g^{11}\left(\frac{\partial g_{21}}{\partial r}
     +\frac{\partial g_{21}}{\partial\theta}-\frac{\partial g_{22}}{\partial r}\right)\\
  &=\frac12\left(-\frac{\partial(\sinh^2r)}{\partial r}\right)=-\frac12\sinh(2r),
\end{align*}
% PDF p. 328
\[
  \Gamma^2_{12}=\Gamma^2_{21},
\]
\[
  \Gamma^2_{12}=\frac12g^{22}\frac{\partial g_{22}}{\partial r}
    =\frac12\frac{1}{\sin^2 r}\,2\sin r\cos r
    =\frac{\cosh r}{\sinh r}=\frac1{\tan r}.
\]
现在用式 \eqref{eq:6A15} 计算协变导数。$\nabla_XY^1$ 为径向分量
\begin{align*}
  \nabla_XY^1&=X^1\frac{\partial Y^1}{\partial r}
      +X^2\left(\frac{\partial Y^1}{\partial\theta}
         -\frac12\sinh(2r)Y^2\right),\\
  \nabla_XY^2&=X^1\left(\frac{\partial Y^2}{\partial r}
        +\frac{1}{\tanh r}Y^2\right)
      +X^2\left(\frac{\partial Y^2}{\partial\theta}
        +\frac{1}{\tanh r}Y^1\right).
\end{align*}
